% Introdução

\section{Introdução}
%\addcontentsline{toc}{chapter}{Introdução} % Adiciona "Introdução" ao sumário
O presente trabalho aborda o desenvolvimento de uma interface web para o Laboratório de Investigação de Filmes Semicondutores (LIFS), com o objetivo de centralizar informações que anteriormente se encontravam dispersas e facilitar o acesso ao conteúdo do laboratório. A proposta surgiu a partir da necessidade de disponibilizar, em um único ambiente, informações sobre o grupo, seus membros, equipamentos, trabalhos acadêmicos e atividades.

O projeto foi desenvolvido considerando principalmente pesquisadores, estudantes e possíveis parceiros externos, buscando uma interface simples, responsiva e acessível. Além da implementação das páginas e funcionalidades previstas, foram realizadas avaliações de desempenho, acessibilidade, responsividade e compatibilidade, permitindo verificar o comportamento da aplicação e identificar pontos que ainda podem ser aprimorados.

Os links e arquivos podem ser obtidos, clicando nos seguintes botões. 

\begin{center}
	
	\noindent
	\begin{minipage}[c]{0.30\textwidth}
		\centering
		\href{https://github.com/rick-12q/SCOM_1_LIFS_SITE}{
			\textcolor{githubcolor}{
				\Large\faGithub
			}
			\\[5pt]
			\textbf{\textcolor{githubcolor}{Repositório}}
		}
	\end{minipage}
	\hfill
	\begin{minipage}[c]{0.30\textwidth}
		\centering
		\href{https://github.com/rick-12q/SCOM_1_LIFS_SITE/blob/main/metricas.zip}{
			\textcolor{metricscolor}{
				\Large\faChartLine
			}
			\\[5pt]
			\textbf{\textcolor{metricscolor}{Métricas}}
		}
	\end{minipage}
	\hfill
	\begin{minipage}[c]{0.30\textwidth}
		\centering
		\href{https://rick-12q.github.io/SCOM_1_LIFS_SITE/}{
			\textcolor{websitecolor}{
				\Large\faGlobe
			}
			\\[5pt]
			\textbf{\textcolor{websitecolor}{Website}}
		}
	\end{minipage}
	
\end{center}

\newpage